\section{Integration, certificate boundaries, and resource behavior}
\label{sec:integration}

The source patch keeps each optional experiment available through a supported
interface. The new production modules have no mandatory third-party Python
dependency. Regina is used for optional geometric checks and the existing
external fallback, not for either the tensor implementation or the cocycle
flow solver.

\subsection{Using the exact tensor backend}

From the \code{fast} directory, the standalone command computes a full
polynomial and reports its one-sided recognition consequence:
\par\noindent\begin{minipage}{\linewidth}
\begin{lstlisting}[language=bash]
python -B -m fastunknot jones examples/trefoil.json --backend tensor
python -B -m fastunknot jones examples/hard_unknot_8.json --backend tensor
\end{lstlisting}
\end{minipage}\par

The first input has a nonidentity polynomial; the second has identity and
the standalone Jones command is inconclusive about knot type. To use the
tensor stage within the full recognizer:
\par\noindent\begin{minipage}{\linewidth}
\begin{lstlisting}[language=bash]
python -B -m fastunknot recognize examples/hard_unknot_8.json \
  --jones-backend tensor
\end{lstlisting}
\end{minipage}\par

Other earlier stages may already decide this small input. Tests explicitly
disable those stages when checking that tensor identity continues to the
Khovanov fallback, and that a separately exhausted fallback returns
\status{UNKNOWN} rather than misusing Jones identity.

The low-level interface is:
\par\noindent\begin{minipage}{\linewidth}
\begin{lstlisting}[language=Python]
from fastunknot import Diagram
from fastunknot.tensor_jones import tensor_jones

diagram = Diagram.from_braid(3, [1, -2] * 5)
result = tensor_jones(
    diagram, certified=True, arithmetic="integer",
    max_states=4096, max_transitions=200000,
)
\end{lstlisting}
\end{minipage}\par

The result contains the exact full polynomial as signed hexadecimal integer
coefficients indexed by integral $t$-exponent, the chosen order and separator
certificate, crossing count, writhe, and measured resource counters. Hex
encoding avoids dependence on a runtime's decimal-digit conversion limit
for very large exact integers. Callers can select \code{"laurent"} for the
reference calculation or \code{"integer-global"} for the arithmetic
ablation. The latter intentionally preserves an inefficient monomial-shift
scheme for comparison.

The retained state and transition limits count mathematical work objects.
They do not independently cap the bytes occupied by a large integer.
Deadline callbacks are checked during preparation, state updates,
coefficient loops, and decoding. A running big-integer operation is not
preempted at an arbitrary machine instruction, so these are cooperative
allowances rather than hard real-time guarantees.

\subsection{Optimizing and checking a cocycle seed}

The numeric API accepts four global vertex IDs and four integral heights
per tetrahedron:
\par\noindent\begin{minipage}{\linewidth}
\begin{lstlisting}[language=Python]
from fastunknot.cocycle_seed import (
    minimize_disc_seed, check_seed_certificate,
)

certificate = minimize_disc_seed(vertices, heights)
checked = check_seed_certificate(vertices, heights, certificate)
assert checked["primal"] == checked["dual"]
\end{lstlisting}
\end{minipage}\par

It rejects malformed rows and nonintegral values. The checker recomputes
the potential-adjusted local gaps, disc counts, shared-vertex matching,
source and target permutations, and objective equality. It does not call
the optimizer or trust its reported statistics.

There are two geometric adapters. The face-pairing API validates reciprocal
gluings, local integer height transitions, and normal matching. The
triangulation-protocol API accepts a separately validated finite
triangulation and integral cocycle, including the archive02 object used in
the tests. The optimization theorem needs only numeric data; the topological
theorems additionally need a genuine compact orientable manifold,
coorientation, and the stated integral class conditions. No adapter silently
infers manifold validity from an optimal flow.

A normal-coordinate vector together with a primal/dual matching certifies
the restricted minimum-disc problem. It does not certify that the surface
is incompressible or that a subsequent hierarchy is essential. Those are
separate certificate types. The proposed Euler-characteristic positive test
in Section~\ref{sec:cocycle} would also require independently bound
knot-exterior provenance and peripheral data before it could decide an
input diagram.

\subsection{Enabling the overlap witness policy}

The new CLI flag enables its required compressed search and relator moves:
\par\noindent\begin{minipage}{\linewidth}
\begin{lstlisting}[language=bash]
python -B -m fastunknot recognize examples/hard_unknot_8.json \
  --group-overlap-witness --group-seconds 2
\end{lstlisting}
\end{minipage}\par

Python callers select the related flags explicitly:
\par\noindent\begin{minipage}{\linewidth}
\begin{lstlisting}[language=Python]
result = recognize(
    diagram,
    use_group=True,
    group_compressed_search=True,
    group_relators=True,
    group_overlap_witness=True,
)
\end{lstlisting}
\end{minipage}\par

The standalone compressed producer accepts
\code{relator\_moves=True, overlap\_witness=True}; the local helper accepts
\code{witness\_first=True}. Non-Boolean values and inconsistent option
combinations are rejected. The adaptive explicit-to-compressed policy is a
different interface and is not silently changed by this option.

Search still has the same resource semantics: a verified successful group
certificate proves the existing conclusion, while an unsuccessful or capped
search remains inconclusive. The first-witness policy may change a future
search trajectory, which is why it is not made the default. Neither
independent verifier needs to know which search policy discovered a move.

\subsection{Repairing large requests to the external normal-surface worker}

The maintained test inventory exposed a separate, reproducible compatibility
bug in the existing worker transport. The parent initially called
\code{communicate(payload, timeout=...)} and retried with no input after a
short timeout. In the tested Python 3.12.14 runtime, the communication code
retained its internal input buffer but registered writable stdin only when
the current call supplied a truthy input argument. A delayed reader with a
500-kilobyte request could therefore receive only the first pipeful, while
the parent waited until its three-second test allowance expired.

The repair gives one communication thread sole ownership of
\code{communicate(payload)}. The caller waits on a completion event in
bounded intervals, checking the existing local deadline and global
cancellation callback on the calling thread. Completion publishes either
the returned streams or the communication error. A worker pipe failure
remains inconclusive, while a caller's cancellation exception propagates.

Cleanup covers every operation after successful child creation, including
event construction, thread construction and startup. It kills a live
child, reaps it, joins a successfully started communicator, and then closes
the pipes. Independent review identified and closed a constructor-scope
gap before the final tests. The scope is the controlled direct-child
Regina worker; it is not a general supervisor for arbitrary grandchildren
retaining inherited pipes.

The existing delayed-reader regression now transmits and recovers all
500,000 input characters within its unchanged three-second allowance.
Existing cancellation and real-engine checks also pass. One additional
test injects communication and thread-setup failures, verifies child
reaping and closed pipes, checks that the payload is submitted once, and
confirms caller-thread polling. The final nine-test worker module passes
in 2.732 seconds, followed by six public-integration tests in 0.010 seconds.
The preserved diagnosis and logs are separate from the performance
benchmarks; this transport repair contributes no claimed mathematical
speedup.


This is an engineering correctness and performance repair, not a new
topological theorem. The external engine remains explicitly trusted for
its own topology operations. Digest validation, response schema checks,
decision semantics, and the complete built-in fallback remain in force.

\subsection{What an independent verifier does and does not establish}

\begin{table}[H]
\centering\small
\begin{tabularx}{\linewidth}{@{}P{0.22\linewidth}YY@{}}
\toprule
Evidence & Independent check & Conclusion\\
\midrule
Separator hierarchy & Graph partition, balance, order, cut width &
The stated width bound for that input and order.\\[3pt]
Full tensor polynomial & State-sum and Regina comparison in the audit &
Implementation agreement on tested inputs; the general exactness theorem
comes from Section~\ref{sec:tensor}.\\[3pt]
Seed potential and matching & Recomputed local spans and equal dual bound &
Global optimality within the vertex-potential family.\\[3pt]
Relator trace & Literal or compressed move-by-move replay &
The existing group certificate conclusion for the original diagram.\\[3pt]
External normal decision & Input digest, process status, schema, engine label &
A completed externally trusted decision, not independent replay of a
geometric hierarchy.\\
\bottomrule
\end{tabularx}
\caption{Trust boundaries retained by the integration. A consistency check
is not automatically a certificate of a stronger mathematical claim.}
\end{table}

All complete output claims refer to a validated input. An invalid planar
encoding, unmatched normal coordinates, an inexact division, or a malformed
certificate is a failed operation. A resource exhaustion is a different
outcome and remains inconclusive. Preserving this distinction is essential
when combining several exact and heuristic discovery stages in one
recognizer.
